\chapter{The Android top-up app}
\label{ch:android}

\chapterbrief
  {Put credit \emph{onto} the card from a phone. Build it yourself, in Dart, from the design you already understand.}
  {Porting a design across languages --- and discovering how much of your stack was incidental.}
  {A phone that adds credit, and cross-language proof your crypto is right.}

\section{Why this chapter is different}

You have built the whole reader stack once. The point of this chapter is not to
learn NFC again. It is to discover something valuable:

\begin{center}\bfseries\itshape
Six of your ten chapters do not exist on Android.
\end{center}

The phone has an NFC controller with a driver already written. Android's
\code{IsoDep} class hands you \code{transceive(byte[]) -> byte[]} directly. That
single method is everything you built in Chapters~\ref{ch:spi}--\ref{ch:apdu}.

\begin{center}\small
\begin{tabular}{@{}L{0.42\linewidth}L{0.48\linewidth}@{}}
\toprule
\thead{Your Pi stack} & \thead{Android} \\
\midrule
Ch.~\ref{ch:spi} \code{SpiLink}         & \emph{gone} --- the OS owns the radio \\
Ch.~\ref{ch:driver} \code{Rc522Driver}  & \emph{gone} \\
Ch.~\ref{ch:reader} \code{Rc522Reader}  & \emph{gone} \\
Ch.~\ref{ch:blocks} \code{ProtocolBlock}& \emph{gone} --- the OS frames I-blocks \\
Ch.~\ref{ch:activation} activation      & \emph{gone} --- you are handed an activated tag \\
Ch.~\ref{ch:apdu} \code{IsoDepChannel}  & \textbf{\code{IsoDep.transceive()}} \\
\midrule
\textbf{Ch.~\ref{ch:cmac}--\ref{ch:secmsg} crypto} & \textbf{you port all of it, unchanged} \\
\bottomrule
\end{tabular}
\end{center}

\begin{concept}{The transport was replaceable; the protocol was the work}
Look at the shape of your \code{ApduChannel} interface from
Chapter~\ref{ch:apdu} and Android's \code{IsoDep.transceive}. They are the same.

That is not luck --- it is what happens when you draw a seam at the right place.
Everything above \code{ApduChannel} ports across a language, an operating system,
and a completely different radio, without changing.

If you want one test of whether an abstraction was well chosen, this is it:
\emph{how much survives when the thing underneath is replaced entirely?}
\end{concept}

\begin{analogy}{Posting the letter}
You have been writing a letter and hand-delivering it. Now you are posting it.
The envelope, the courier, the van --- all different. The letter is identical.
\end{analogy}

\section{Setup}

\begin{shell}
flutter create --org com.midnightbrewer --platforms android topup_app
cd topup_app
\end{shell}

\begin{xcode}
dependencies:
  flutter:
    sdk: flutter
  # Android IsoDep (ISO 14443-4) transceive.
  # 4.x has a different API from 3.x -- pin the major version.
  nfc_manager: ^4.2.1
  # AES-ECB/CBC and AES-CMAC.
  pointycastle: ^4.0.0
\end{xcode}

And in \code{android/app/src/main/AndroidManifest.xml}:

\begin{xcode}
<uses-permission android:name="android.permission.NFC" />
<uses-feature android:name="android.hardware.nfc" android:required="true" />
\end{xcode}

\section{Dart for a Java programmer}

The languages are close cousins. Four differences matter here.

\begin{description}[leftmargin=0pt, style=nextline]
  \item[\code{Uint8List}, not \code{byte[]}.] From \code{dart:typed\_data}.
        Crucially its elements are \textbf{unsigned 0--255}: Dart has no
        signed-byte trap, and all those \code{\& 0xFF} masks simply vanish.
  \item[\code{int} is 64-bit.] So \code{1 << 40} works fine, but when writing a
        32-bit value you must mask deliberately: \code{(v >> 24) \& 0xFF}.
  \item[Async is explicit.] NFC calls are \code{Future}s. A function that awaits
        must be \code{async}, and its callers must \code{await} it --- the
        colouring propagates upward, much as \code{throws SpiException} did in
        Java. Same discipline, different keyword.
  \item[No checked exceptions.] Dart has \code{throw}/\code{catch} but the
        compiler forces nothing. You lose the nagging, so you must be more
        deliberate about error paths.
\end{description}

\section{Build it}

\subsection{The transport --- one class replaces six chapters}

\begin{dcode}
import 'dart:typed_data';
import 'package:nfc_manager/nfc_manager.dart';

/// Sends one APDU to an activated card and returns the reply.
///
/// This is the ENTIRE transport layer -- the Android NFC stack has already
/// done activation, anticollision, RATS and ISO-DEP framing for us.
/// Compare Chapters 2-7 of the Pi build. Same seam, no work.
class IsoDepTransceiver {
  IsoDepTransceiver(this._isoDep);

  final IsoDep _isoDep;

  Future<Uint8List> transceive(Uint8List apdu) async {
    // TODO: return await _isoDep.transceive(data: apdu);
    //   Wrap failures in your own exception type so the UI never sees a
    //   platform exception -- the same translate() discipline as Chapter 12.
    throw UnimplementedError();
  }
}
\end{dcode}

\subsection{The crypto --- port it directly}

Your Java from Chapters~\ref{ch:cmac}--\ref{ch:secmsg} translates almost line for
line. \code{pointycastle} supplies AES \emph{and} CMAC, so you do not even
reimplement RFC 4493 --- but you must verify it produces the same bytes.

\begin{dcode}
import 'dart:typed_data';
import 'package:pointycastle/export.dart';

const int kAesBlockSize = 16;

Uint8List aesEcbEncryptBlock(Uint8List key, Uint8List block) {
  final engine = AESEngine()..init(true, KeyParameter(key));
  final out = Uint8List(kAesBlockSize);
  engine.processBlock(block, 0, out, 0);
  return out;
}

Uint8List aesCbcEncrypt(Uint8List key, Uint8List iv, Uint8List data) {
  // TODO: CBCBlockCipher(AESEngine())
  //         ..init(true, ParametersWithIV(KeyParameter(key), iv));
  //       process block by block; NO padding (DESFire pads itself).
  throw UnimplementedError();
}

Uint8List aesCmac(Uint8List key, Uint8List message) {
  final mac = CMac(AESEngine(), 128)..init(KeyParameter(key));
  return mac.process(message);
}

/// DESFire keeps the ODD-indexed bytes: mac[1], mac[3], ... mac[15].
Uint8List truncateMac(Uint8List mac) {
  // TODO: 8 bytes, out[i] = mac[i * 2 + 1];
  throw UnimplementedError();
}

/// The 32-byte derivation vector, exactly as in Chapter 9.
Uint8List deriveSessionKey(
    Uint8List key, Uint8List rndA, Uint8List rndB, int label) {
  // TODO: sv[0] = label (0xA5 enc / 0x5A mac); sv[1] = the other;
  //       sv[2..5] = 00 01 00 80;
  //       sv[6..7] = rndA[0..1];
  //       sv[8..13] = rndA[2..7] ^ rndB[0..5];
  //       sv[14..23] = rndB[6..15];
  //       sv[24..31] = rndA[8..15];
  //       return aesCmac(key, sv);
  throw UnimplementedError();
}

/// ISO 7816-4: append 0x80 then zeros. ALWAYS adds at least one byte.
Uint8List padIso7816(Uint8List data) {
  final paddedLength = (data.length | 0x0F) + 1;
  final out = Uint8List(paddedLength);
  out.setRange(0, data.length, data);
  out[data.length] = 0x80;
  return out;
}
\end{dcode}

\begin{verifybox}
\textbf{Verify with the same vectors.} This is the elegant part: the vectors in
Appendix~\ref{app:vectors} are language-agnostic. Your Dart must produce
byte-for-byte what your Java produced.

\begin{dcode}
void main() {
  test('derives the Chapter 9 session keys', () {
    final key  = fromHex('00112233445566778899aabbccddeeff');
    final rndA = fromHex('a0a1a2a3a4a5a6a7a8a9aaabacadaeaf');
    final rndB = fromHex('b0b1b2b3b4b5b6b7b8b9babbbcbdbebf');

    expect(toHex(deriveSessionKey(key, rndA, rndB, 0xA5)),
           'd46c03528ecd3b1f1148e708d3e9729d');
    expect(toHex(deriveSessionKey(key, rndA, rndB, 0x5A)),
           '116d156cc502da820eb57915f189ec74');
  });

  test('truncates a CMAC to the odd bytes', () {
    expect(toHex(truncateMac(fromHex('00112233445566778899aabbccddeeff'))),
           '1133557799bbddff');
  });
}
\end{dcode}

\begin{shell}
flutter test           # no phone, no card -- same discipline as the Pi build
\end{shell}

If your Dart disagrees with your Java on any vector, \textbf{one of them is
wrong}, and you now know before either touches a card. Two independent
implementations converging on the same bytes is much stronger evidence than
either alone.
\end{verifybox}

\subsection{Credit --- the mirror image of debit}

\begin{dcode}
/// Credit (0x0C) in CommMode.Full -- the mirror image of Debit.
///
/// Identical wire format; only the command byte differs. Like debit it only
/// STAGES the change: nothing moves until commitTransaction().
Future<void> credit(int fileNo, int amountCents) async {
  // TODO: same shape as Chapter 11's debit:
  //   header     = [fileNo]
  //   ciphertext = session.encryptCommandData(encodeValue(amountCents))
  //   macInput   = header + ciphertext
  //   mac        = session.commandMac(0x0C, macInput)
  //   send, check status, incrementCounter, verify response MAC over empty
  throw UnimplementedError();
}
\end{dcode}

And the top-up sequence, which reads the balance back:

\begin{dcode}
Future<int> topUp(int fileNo, int amountCents) async {
  final before = await getValue(fileNo);
  await credit(fileNo, amountCents);
  await commitTransaction();
  final after = await getValue(fileNo);      // trust nothing; re-read

  if (after != before + amountCents) {
    throw StateError('balance did not move as expected');
  }
  return after;
}
\end{dcode}

Reading the balance back after the commit is not paranoia --- it is the only way
to \emph{know}. If the phone loses the field between \code{credit} and
\code{commitTransaction}, the card discards the delta and the balance is
unchanged. Same safety property as Chapter~\ref{ch:secmsg}, now protecting your
customer's money rather than yours.

\section{Provisioning: turning a blank card into a wallet}

A factory-fresh DESFire has no application, no value file, and a DES master key.
To make one work with your kiosk:

\begin{enumerate}
  \item authenticate at PICC level with the master key (16 $\times$ \hex{0x00} factory default);
  \item \code{CreateApplication} (\hex{0xCA}) --- AID \hex{0x010203}, 3 keys, AES;
  \item select the application, authenticate with its master key (initially 16 $\times$ \hex{0x00});
  \item \code{ChangeKey} (\hex{0xC4}) --- set app master key (\hex{0x11}\ldots) and user key (\hex{0x22}\ldots);
  \item re-authenticate with the new app master key;
  \item \code{CreateValueFile} (\hex{0xCC}) --- file \hex{0x01}, CommMode.Full,
        access via key 2, limits 0\ldots100000, initial value 0.
\end{enumerate}

Wire format for \code{CreateValueFile}, 17 bytes:

\begin{trace}
[0]      fileNo
[1]      commSettings          0x03 = Full
[2..3]   accessRights          little-endian; nibbles: read|write|rw|changeRights
[4..7]   lowerLimit            signed LE
[8..11]  upperLimit            signed LE
[12..15] initialValue          signed LE
[16]     limitedCreditEnabled  0x00 / 0x01
\end{trace}

The access-rights word packs four 4-bit key numbers: \emph{read} (GetValue),
\emph{write} (Debit), \emph{readWrite} (Credit too), and
\emph{changeAccessRights}. Using key 2 for all four gives \hex{0x2222}. Nibble
\hex{0xE} means free access; \hex{0xF} means never.

\begin{trap}{\texttt{ChangeKey} is irreversible}
Set a key to a value you do not record and that card is scrap --- permanently, by
design. There is no recovery path; that \emph{is} the security property.

Before you run provisioning:
\begin{itemize}
  \item use \textbf{dedicated test cards} you bought for this project;
  \item never change the PICC master key from its factory default (this project
        deliberately does not, so any other tool can re-personalise the card);
  \item never point this at a card you did not buy for this --- the DESFire in
        your wallet may be your office badge or transit pass, and it will not
        survive.
\end{itemize}

And \textbf{never auto-provision on tap}. Require an explicit ``Set up this
card'' button press. A member of the public tapping your app must never have
their existing card silently rewritten.
\end{trap}

\section{The security reality of hardcoded keys}
\label{sec:realkeys}

Your app will contain, in plaintext, the AES keys for every card in the fleet.
Anyone who pulls the APK off a phone can extract them and mint unlimited balance
on every card you have issued.

\textbf{For a bench project with test cards this is fine.} Say so out loud, in a
comment, so nobody mistakes it for a design:

\begin{dcode}
/// BENCH KEYS. These are compiled into the APK in plaintext. Anyone who pulls
/// the APK can read them and mint unlimited balance on every card, because the
/// same key is on every card. Acceptable for a bench rig with test cards.
/// NOT acceptable for anything representing real money.
\end{dcode}

\begin{concept}{What a real deployment does instead}
Worth understanding even if you never build it.

\textbf{1. Key diversification.} Each card gets a \emph{different} key, derived
from the card's UID and a master key (typically via CMAC). Compromise one card
and you learn one card's key --- not the fleet's.

\textbf{2. A backend HSM or SAM.} The master key never leaves tamper-resistant
hardware. The phone becomes a dumb pipe: it relays APDUs between the card and the
server, and the \emph{server} performs the cryptography.

\textbf{3. Server-side transaction logs}, so the balance on the card is a cache
rather than the source of truth, and double-spends are detectable after the fact.

The gap between ``works'' and ``safe to deploy'' is mostly key management, and it
is why payment systems have HSMs. Knowing exactly where your bench rig sits on
that spectrum is the professional part.
\end{concept}

\section{Build order}

Do it in this order:

\begin{enumerate}
  \item \code{flutter create} plus dependencies; blank app runs on your phone.
  \item \textbf{Crypto module and tests.} No NFC yet. All vectors green.
  \item \code{IsoDepTransceiver} and a ``read the card's version'' screen. Prove
        \hex{90 60 00 00 00} returns \hex{...91AF}.
  \item \code{authenticateEv2First}, on a card the kiosk already works with.
  \item \code{getValue} --- display the balance. Read-only, so mistakes are harmless.
  \item \code{credit} $+$ \code{commit} $+$ read-back. Now money moves.
  \item Provisioning screen, last, behind an explicit button.
  \item UI polish --- the NFC pulse animation, preset amounts, result screen.
\end{enumerate}

Steps 2 and 5 are the natural stopping points if you are filming: ``the crypto
agrees across two languages'' and ``the phone reads the balance the Pi wrote''.

\section{Traps}

\begin{traptable}
\code{nfc\_manager} API does not match tutorials & 3.x and 4.x differ. Pin 4.x and use its docs. \\
Works on one phone, not another & NFC antenna placement varies enormously. Try the upper third of the back. \\
Tag lost mid-transaction & Phone moved. This is \emph{normal} --- it is why commit exists. Handle it. \\
Dart and Java disagree on a MAC & Almost always the counter or TI, not the AES. Print both inputs. \\
\hex{91 AE} after provisioning & You changed a key and re-authenticated with the old one. \\
Card bricked & You changed a key you did not record. There is no recovery. This is why you buy spares. \\
\end{traptable}

\section{Commit}

\begin{shell}
git add -A
git commit -m "topup: Flutter app -- crypto ported, credit + commit, provisioning"
\end{shell}
